% !TEX root = ../../hw05.tex
\begin{proof}[Solution]
  To disprove uniform continuity, we need one fixed output tolerance
  that no single input tolerance can guarantee over all of $\R$.
  Choose $\varepsilon_0=1$.

  Let $\delta>0$ be arbitrary. By the Archimedean property, choose
  $n\in\N$ with $n>1/\delta$, and set
  \[
    x_n=n,\qquad y_n=n+\frac1n.
  \]
  These inputs are less than $\delta$ apart, since
  \[
    |y_n-x_n|=\frac1n<\delta.
  \]
  However, their outputs satisfy
  \begin{align*}
    |f(y_n)-f(x_n)|
      &=\left(n+\frac1n\right)^2-n^2\\
      &=n^2+2+\frac1{n^2}-n^2\\
      &=2+\frac1{n^2}>1=\varepsilon_0.
  \end{align*}
  Thus every $\delta>0$ fails for the same $\varepsilon_0=1$, so $f$
  is not uniformly continuous on $\R$.

  What goes wrong is visible in the factorisation
  \[
    |y^2-x^2|=|y-x||y+x|.
  \]
  Even when $|y-x|$ is small, the factor $|y+x|$ can be arbitrarily
  large. In our construction the inputs get closer together as they
  move to the right, while the output gap stays above $2$.
  \autoref{fig:hw05:square-output-gap} shows one pair from this sequence.

  \begin{marginfigure}
    \centering
    % !TEX root = ../../problems/hw05.tex
% A magnified local view of x^2 at x_n=3 and y_n=3+1/3.
% Axis origins lie outside this view; horizontal and vertical scales differ.
% Shared physical dimensions for interval and number-line diagrams.
\tikzset{
  interval diagram/.style={
    x=1cm,y=6mm,font=\footnotesize,
    color=black,line width=0.5pt,line cap=round,line join=round,
    >={Stealth[length=4pt,width=3pt]}
  },
  interval endpoint/.style={
    circle,draw=black,line width=0.5pt,
    minimum size=4pt,inner sep=0pt,outer sep=0pt
  },
  interval open/.style={interval endpoint,fill=white},
  interval closed/.style={interval endpoint,fill=black},
  interval label/.style={inner sep=1pt}
}

\begin{tikzpicture}[interval diagram,x=60mm,y=9mm]
  \draw[->] (2.9,8.6) -- (3.5,8.6) node[right=4pt] {$x$};
  \draw[->] (2.9,8.6) -- (2.9,12.2) node[above=4pt] {$y$};
  \draw[<->,domain=2.98:3.43,samples=80] plot (\x,{\x*\x});

  \draw[gray,densely dashed] (3,8.6) -- (3,9) -- (3.333333,9);
  \draw[gray,densely dashed] (3.333333,8.6) -- (3.333333,11.111111);
  \draw[gray] (3.333333,9) -- (3.44,9);
  \draw[gray] (3.333333,11.111111) -- (3.44,11.111111);
  \fill (3,9) circle[radius=1.6pt];
  \fill (3.333333,11.111111) circle[radius=1.6pt];
  \node[below=6pt] at (3,8.6) {$x_n$};
  \node[below=6pt] at (3.333333,8.6) {$y_n$};

  \draw[<->] (3,7.3) -- (3.333333,7.3)
    node[midway,below=4pt] {$1/n$};
  \draw[<->] (3.44,9) -- (3.44,11.111111)
    node[midway,right=4pt] {$>2$};
\end{tikzpicture}

    \caption[Small input gaps with large output gaps.]{%
      A local view of $y=x^2$ for the pair $x_n=n$,
      $y_n=n+1/n$ when $n=3$. A horizontal gap of $1/n$ produces a
      vertical gap of $2+1/n^2$. The two axes use different scales.}
    \label{fig:hw05:square-output-gap}
  \end{marginfigure}
\end{proof}
